\section*{Bab \ref{ch:g8:fractions} --- Pecahan: Keempat Operasinya}

\begin{solution}{exo:g8:fractions:1}
$\dfrac34 + \dfrac25 = \dfrac{15}{20} + \dfrac{8}{20} =
\dfrac{23}{20}$.

$\dfrac56 - \dfrac38 = \dfrac{20}{24} - \dfrac{9}{24} =
\dfrac{11}{24}$.

$\dfrac{7}{10} + \dfrac{8}{15} = \dfrac{21}{30} + \dfrac{16}{30} =
\dfrac{37}{30}$.
\end{solution}

\begin{solution}{exo:g8:fractions:2}
$-\dfrac27 + \dfrac{3}{14} = -\dfrac{4}{14} + \dfrac{3}{14} =
-\dfrac{1}{14}$.

$\dfrac16 - \dfrac54 = \dfrac{2}{12} - \dfrac{15}{12} =
-\dfrac{13}{12}$.

$-\dfrac35 - \dfrac12 = -\dfrac{6}{10} - \dfrac{5}{10} =
-\dfrac{11}{10}$.
\end{solution}

\begin{solution}{exo:g8:fractions:3}
$\dfrac58 \times \dfrac{4}{15} = \dfrac{5 \times 4}{8 \times 15}
= \dfrac{1}{2 \times 3} = \dfrac16$ (sederhanakan dengan $5$ dan dengan $4$).

$\dfrac{9}{14} \times \dfrac73 = \dfrac{9 \times 7}{14 \times 3}
= \dfrac{3}{2}$ (sederhanakan dengan $7$ dan dengan $3$).

$\left(-\dfrac65\right) \times \dfrac{10}{9} = -\dfrac{60}{45}
= -\dfrac43$ (tandanya berlawanan: negatif).
\end{solution}

\begin{solution}{exo:g8:fractions:4}
\omterm{def:g8:fractions:inverse}{Kebalikannya}: $\dfrac73$; \quad $\dfrac15$; \quad $4$; \quad
$-\dfrac92$ (\omterm{def:g8:fractions:inverse}{kebalikannya} mempertahankan tandanya).
\end{solution}

\begin{solution}{exo:g8:fractions:5}
$\dfrac49 \div \dfrac23 = \dfrac49 \times \dfrac32 = \dfrac{12}{18}
= \dfrac23$.

$\dfrac75 \div 14 = \dfrac75 \times \dfrac{1}{14} = \dfrac{7}{70}
= \dfrac{1}{10}$.

$6 \div \dfrac34 = 6 \times \dfrac43 = 8$.
\end{solution}

\begin{solution}{exo:g8:fractions:6}
$\dfrac{3/4}{9/8} = \dfrac34 \times \dfrac89 = \dfrac{24}{36} =
\dfrac23$.

$\dfrac{5/6}{10} = \dfrac56 \times \dfrac{1}{10} = \dfrac{5}{60} =
\dfrac{1}{12}$.

$\dfrac{2}{4/7} = 2 \times \dfrac74 = \dfrac72$.
\end{solution}

\begin{solution}{exo:g8:fractions:7}
\omterm{def:g2:mult:def}{Hasil kalinya} dahulu:
$\dfrac13 + \dfrac12 \times \dfrac45 = \dfrac13 + \dfrac{4}{10}
= \dfrac{10}{30} + \dfrac{12}{30} = \dfrac{22}{30} = \dfrac{11}{15}$.

Tanda kurungnya dahulu:
$\left(\dfrac13 + \dfrac12\right) \times \dfrac45
= \dfrac56 \times \dfrac45 = \dfrac{20}{30} = \dfrac23$.
\end{solution}

\begin{solution}{exo:g8:fractions:8}
Terisi: $\dfrac25 + \dfrac13 = \dfrac{6}{15} + \dfrac{5}{15} =
\dfrac{11}{15}$ bagian tangkinya. Kosong:
$1 - \dfrac{11}{15} = \dfrac{4}{15}$.
\end{solution}

\begin{solution}{exo:g8:fractions:9}
Bagian yang memperoleh lebih dari $14$:
$\dfrac23 \times \dfrac34 = \dfrac{6}{12} = \dfrac12$ dari kelasnya.
Kalau itu $12$ murid, kelasnya berisi $24$ murid.
\end{solution}

\begin{solution}{exo:g8:fractions:10}
$\dfrac{15}{2} \div \dfrac34 = \dfrac{15}{2} \times \dfrac43
= \dfrac{60}{6} = 10$ potongan --- bilangan bulat, tanpa tali yang terbuang.
\end{solution}

\begin{solution}{exo:g8:fractions:11}
$E$: pembagiannya dahulu:
$\dfrac56 \div \dfrac{10}{9} = \dfrac56 \times \dfrac{9}{10}
= \dfrac{45}{60} = \dfrac34$; lalu
$E = \dfrac23 - \dfrac34 = \dfrac{8}{12} - \dfrac{9}{12} =
-\dfrac{1}{12}$.

$F = \left(-\dfrac12\right)^2 \times \dfrac83 = \dfrac14 \times
\dfrac83 = \dfrac{8}{12} = \dfrac23$ (kuadratnya positif).
\end{solution}

\begin{solution}{exo:g8:fractions:12}
Dari dalam ke luar: $1 + 1 = 2$; lalu
$1 + \dfrac12 = \dfrac32$; lalu
\[
1 + \frac{1}{3/2} = 1 + \frac23 = \frac53 .
\]
\end{solution}

\begin{solution}{pb:g8:fractions:1}
\textbf{1.} $\frac12 + \frac13 = \frac36 + \frac26 = \frac56$ dan
$\frac12 + \frac14 = \frac24 + \frac14 = \frac34$. Jadi
$\frac56 = \frac12 + \frac13$ dan $\frac34 = \frac12 + \frac14$
adalah tulisan Mesir (\omterm{def:g6:fractions:def}{pecahan} satuannya berbeda pada masing-masing).

\textbf{2.} Dengan \omterm{def:g4:fractions:def}{penyebut} yang sama $28$:
\[
\frac14 + \frac{1}{28} = \frac{7}{28} + \frac{1}{28}
= \frac{8}{28} = \frac{2}{7} .
\]

\textbf{3.} $\frac13 + \frac16 = \frac26 + \frac16 = \frac36 =
\frac12$, dan $\frac14 + \frac{1}{12} = \frac{3}{12} +
\frac{1}{12} = \frac{4}{12} = \frac13$. Secara umum, \omterm{def:g4:fractions:def}{penyebut}
bersama $\frac{1}{n+1}$ dan $\frac{1}{n(n+1)}$ adalah
$n(n+1)$:
\[
\frac{1}{n+1} + \frac{1}{n(n+1)}
= \frac{n}{n(n+1)} + \frac{1}{n(n+1)}
= \frac{n+1}{n(n+1)}
= \frac{1}{n} .
\]

\textbf{4.} Mulailah dari $1 = \frac12 + \frac12$ lalu pecahlah
setengah yang kedua menurut pertanyaan 3:
\[
1 = \frac12 + \frac13 + \frac16 ,
\]
yaitu tiga \omterm{def:g6:fractions:def}{pecahan} satuan yang berbeda-beda.

\textbf{5.} Memecah $\frac13$ pada $\frac56 = \frac12 +
\frac13$ memberi
\[
\frac56 = \frac12 + \frac14 + \frac{1}{12} .
\]
Tidak ada tulisan Mesir yang pernah tunggal: kesamaan pemecahannya
selalu dapat mengganti \omterm{def:g6:fractions:def}{pecahan} satuan yang \omterm{def:g4:fractions:def}{penyebutnya} terbesar
dengan dua \omterm{def:g6:fractions:def}{pecahan} satuan baru yang bahkan lebih kecil --- jadi dari
tulisan mana pun kita dapat membuat tulisan yang lain.

\textbf{6.} Dengan \omterm{def:g4:fractions:def}{penyebut} $21$: $\frac27 = \frac{6}{21} <
\frac{7}{21} = \frac13$. Dengan \omterm{def:g4:fractions:def}{penyebut} $28$: $\frac14 =
\frac{7}{28} < \frac{8}{28} = \frac27$. \omterm{def:g6:fractions:def}{Pecahan} satuan mengecil
ketika \omterm{def:g4:fractions:def}{penyebutnya} membesar, $1 > \frac12 > \frac13 > \frac14 >
\dots$; karena $\frac13$ (dan semua yang di atasnya) terlalu besar
sedangkan $\frac14$ muat, $\frac14$ adalah \omterm{def:g6:fractions:def}{pecahan} satuan terbesar
yang lebih kecil daripada $\frac27$.

\textbf{7.} $\frac27 - \frac14 = \frac{8}{28} - \frac{7}{28} =
\frac{1}{28}$: cara rakusnya menghasilkan
$\frac27 = \frac14 + \frac{1}{28}$, persis catatan juru tulis pada
pertanyaan 2.

\textbf{8.} Dengan \omterm{def:g4:fractions:def}{penyebut} $14$: $\frac12 = \frac{7}{14} <
\frac{10}{14} = \frac57$, sedangkan satu-satunya \omterm{def:g6:fractions:def}{pecahan} satuan yang
lebih besar, yaitu $1$, melampaui $\frac57$. Jadi $\frac12$ yang
terbesar yang muat, dan
\[
\frac57 - \frac12 = \frac{10}{14} - \frac{7}{14} = \frac{3}{14} .
\]

\textbf{9.} Dengan \omterm{def:g4:fractions:def}{penyebut} $28$: $\frac{3}{14} = \frac{6}{28} <
\frac{7}{28} = \frac14$, jadi $\frac14$ terlalu besar. Dengan
\omterm{def:g4:fractions:def}{penyebut} $70$: $\frac15 = \frac{14}{70} < \frac{15}{70} =
\frac{3}{14}$, jadi $\frac15$ muat: itulah \omterm{def:g6:fractions:def}{pecahan} satuan terbesar
yang lebih kecil daripada $\frac{3}{14}$. Lalu
\[
\frac{3}{14} - \frac15 = \frac{15}{70} - \frac{14}{70}
= \frac{1}{70},
\qquad\text{jadi}\qquad
\frac57 = \frac12 + \frac15 + \frac{1}{70} .
\]
Periksa: $\frac{35}{70} + \frac{14}{70} + \frac{1}{70} =
\frac{50}{70} = \frac57$.

\textbf{10.} Langkah pertama: $\frac12 = \frac{5}{10} < \frac{8}{10} =
\frac45$ dan $1 > \frac45$, jadi kurangkan $\frac12$:
$\frac45 - \frac12 = \frac{8}{10} - \frac{5}{10} = \frac{3}{10}$.
Langkah kedua: $\frac13 = \frac{10}{30} > \frac{9}{30} =
\frac{3}{10}$ terlalu besar, sedangkan $\frac14 = \frac{5}{20} <
\frac{6}{20} = \frac{3}{10}$ muat, jadi kurangkan $\frac14$:
$\frac{3}{10} - \frac14 = \frac{6}{20} - \frac{5}{20} =
\frac{1}{20}$. Kesimpulannya:
\[
\frac45 = \frac12 + \frac14 + \frac{1}{20},
\qquad
\frac{10}{20} + \frac{5}{20} + \frac{1}{20} = \frac{16}{20}
= \frac45 . \checkmark
\]

\textbf{11.} \omterm{def:g4:fractions:def}{Pembilang} \omterm{def:g3:division:remainder}{sisa} yang berturut-turut adalah:
untuk $\frac27$: $2$, lalu $1$; untuk $\frac57$: $5$, lalu $3$, lalu
$1$; untuk $\frac45$: $4$, lalu $3$, lalu $1$. Tiap daftarnya
\emph{turun secara tegas}.

\textbf{12.} \omterm{def:g4:fractions:def}{Penyebut} bersama $\frac ab$ dan $\frac1n$
adalah $b \times n$:
\[
\frac{a}{b} - \frac{1}{n}
= \frac{a \times n}{b \times n} - \frac{b}{b \times n}
= \frac{a \times n - b}{b \times n} .
\]

\textbf{13.} Pada \omterm{def:g4:fractions:def}{penyebut} bersama $b \times (n-1)$, ketaksamaan
$\frac{1}{n-1} > \frac ab$ membandingkan \omterm{def:g4:fractions:def}{pembilang}
$b$ dan $a \times (n - 1)$:
\[
b > a \times (n - 1) = a \times n - a ,
\]
dan menambahkan $a - b$ pada kedua ruasnya memberi
$a > a \times n - b$. Jadi sesudah tiap langkah rakusnya, \omterm{def:g4:fractions:def}{pembilang}
\omterm{def:g3:division:remainder}{sisanya} (pertanyaan 12) adalah bilangan bulat yang benar-benar lebih
kecil daripada sebelumnya. Bilangan bulat tidak dapat turun secara
tegas selamanya: sesudah paling banyak $a$ langkah, \omterm{def:g3:division:remainder}{sisanya}
berpembilang $1$ --- yaitu \omterm{def:g6:fractions:def}{pecahan} satuan --- atau bernilai
$0$, dan caranya berhenti.

\textbf{14.} Dengan \omterm{def:g4:fractions:def}{penyebut} yang sama $18$:
\[
\frac12 + \frac13 + \frac19
= \frac{9}{18} + \frac{6}{18} + \frac{2}{18}
= \frac{17}{18} ,
\]
yang \emph{lebih kecil daripada} $1$: wasiatnya melupakan
$\frac{1}{18}$ kawanannya. Celah itulah seluruh triknya. Dari $18$
ekor unta, bagian $\frac12$, $\frac13$, $\frac19$ adalah bilangan
bulat $9$, $6$ dan $2$, yang berjumlah $9 + 6 + 2 = 17$: persis
kawanannya, dan unta kedelapan belasnya pulang. Dan tidak ada yang
dapat mengeluh: dari $17$ ekor unta, wasiatnya menjanjikan
$\frac{17}{2} = 8.5$, lalu $\frac{17}{3} \approx 5.67$, lalu
$\frac{17}{9} \approx 1.89$ ekor unta, dan tiap anaknya menerima
benar-benar lebih banyak ($9$, $6$ dan $2$). Yang terlupakan,
$\frac{1}{18}$, adalah yang untuk sementara diisi unta tetangganya.

\textbf{15.} Pertanyaan 4 menyediakan wasiat yang sudah diperbaiki:
\[
\frac12 + \frac13 + \frac16 = 1 .
\]
Bagiannya berupa bilangan bulat kalau besar kawanannya \omterm{def:g6:wholes:divisible}{habis dibagi}
$2$, \omterm{def:g6:wholes:divisible}{habis dibagi} $3$ dan \omterm{def:g6:wholes:divisible}{habis dibagi} $6$ --- yaitu \omterm{def:g6:wholes:divisible}{habis dibagi}
$6$. Kawanan terkecil seperti itu adalah $6$ ekor unta: bagiannya
$3$, $2$ dan $1$, dan tidak ada yang tersisa.

\end{solution}
